\exercise{7.3}{A catalogue of twenty sequences and their limits}

\goal Determine convergence or divergence for every sequence in
parts (a)--(t), and give the real limit whenever one exists.
Although the exercise does not require proofs, we include a justification
of every answer. The arguments use the definition of convergence
directly, rather than the limit laws developed later in Section~9.

\idea For a convergent sequence, subtract the proposed limit and
bound the resulting error by a reciprocal that becomes arbitrarily
small. For a divergent sequence, find either arbitrarily late values
that are arbitrarily large, or two arbitrarily late values separated
by a fixed positive distance.

\subsubsection*{Elementary facts used in the individual parts}

\step{Fact 1: Reciprocal error estimates imply convergence.}
The definition says that $u_n\to L$ if, for every $\varepsilon>0$,
there is a threshold $N$ such that
\[
 n>N\quad\Longrightarrow\quad\abs{u_n-L}<\varepsilon.
\]
Suppose we prove $\abs{u_n-L}\leq C/n$ for every $n\geq n_0$,
where $C>0$ is fixed. Given $\varepsilon>0$, choose a positive
integer $N>\max\{n_0,C/\varepsilon\}$, using the Archimedean
property. Then $n>N$ implies both $n\geq n_0$ and
$n>C/\varepsilon$. Multiplying the latter inequality by
$\varepsilon>0$ and dividing by $n>0$ gives
\[
 \abs{u_n-L}\leq\frac Cn<\varepsilon.
\]
This proves convergence from the definition. In particular,
$1/n\to0$. If an error is identically $0$, it is less than every
$\varepsilon>0$ without any additional estimate.

We will also use a shifted version. If, for all sufficiently large
$n$, $\abs{u_n-L}\leq C/(n-r)$, with fixed real $r$ and $n>r$,
choose $N$ larger than the starting index and $r+C/\varepsilon$.
Then $n>N$ gives $n-r>C/\varepsilon$, and the same calculation
gives $\abs{u_n-L}<\varepsilon$. This proves the shifted version
without any new limit theorem.

Only the tail is used in either proof. Thus a finite number of unusual
initial values does not affect the conclusion: increase $N$ to exceed
all their indices, and every $n>N$ still satisfies the required
estimate. This observation will matter when an initial denominator
is negative but never zero.

\step{Fact 2: Powers of two dominate a simple linear expression.}
For every integer $m\geq0$,
\[
 2^m\geq m+1.
\]
At $m=0$, both sides equal $1$. If the inequality holds at $m$,
then multiplication by $2>0$ gives
\[
 2^{m+1}=2\cdot2^m\geq2(m+1)=2m+2\geq m+2,
\]
where the last difference is $(2m+2)-(m+2)=m\geq0$.
Induction proves the assertion. Taking reciprocals of positive
quantities consequently gives
\[
 0<2^{-m}\leq\frac1{m+1}.
\]

\step{Fact 3: Separated arbitrarily late values rule out a real limit.}
Suppose some fixed $D>0$ has this property: for every threshold
$N$, there are indices $p,q>N$ with $\abs{u_p-u_q}\geq D$.
Assume, to obtain a contradiction, that $u_n\to L\in\R$.
Apply the definition with $\varepsilon=D/3$. It supplies $N$ such
that both late terms obey
$\abs{u_p-L}<D/3$ and $\abs{u_q-L}<D/3$. The triangle inequality
would then give
\[
 \abs{u_p-u_q}
 =\abs{(u_p-L)+(L-u_q)}
 \leq\abs{u_p-L}+\abs{L-u_q}
 <\frac D3+\frac D3=\frac{2D}{3}<D,
\]
contradicting the stipulated separation. Thus such a sequence
does not converge to a real number. This argument does not assume
any theorem about subsequences.

\step{Fact 4: Arbitrarily large late values also rule out a real limit.}
Suppose that for every $B\in\R$ and every threshold $N$, some
$n>N$ satisfies $u_n>B$. If $u_n\to L\in\R$, the definition
with $\varepsilon=1$ would supply $N$ for which
$\abs{u_n-L}<1$ whenever $n>N$. Choose a late value with
$u_n>L+1$. It satisfies $u_n-L>1$, and therefore
$\abs{u_n-L}=u_n-L>1$, a contradiction. This is a direct
nonconvergence argument; it does not invoke a later boundedness
theorem for convergent sequences.

\subsubsection*{Solutions of parts (a)--(t)}

\step{(a) The sequence $a_n=n/(n+1)$.}
For every $n\in\N$, $n+1\geq2>0$, so every term is defined.
To test the candidate $1$, subtract it using the common denominator
$n+1$:
\[
 a_n-1=\frac n{n+1}-\frac{n+1}{n+1}
 =\frac{n-(n+1)}{n+1}
 =-\frac1{n+1}.
\]
It follows that
\[
 \abs{a_n-1}=\frac1{n+1}\leq\frac1n.
\]
The inequality uses $n+1\geq n>0$ and the reversal of order for
positive reciprocals. Fact 1 with $C=1$ now proves the limit.

\answer{Part (a) converges, and $\displaystyle\lim_{n\to\infty}\frac n{n+1}=1$.}

\step{(b) The sequence $b_n=(n^2+3)/(n^2-3)$.}
The denominator is nonzero at every positive integer. At $n=1$
it is $1-3=-2$, while $n\geq2$ implies $n^2\geq4$ and
$n^2-3\geq1>0$. Thus there is no undefined term, including in
the initial part of the sequence.

Subtracting the candidate $1$ gives
\[
 b_n-1
 =\frac{n^2+3-(n^2-3)}{n^2-3}
 =\frac{n^2+3-n^2+3}{n^2-3}
 =\frac6{n^2-3}.
\]
For $n\geq3$, we have $n^2\geq9\geq6$, so $3\leq n^2/2$.
Subtracting $3$ from $n^2$ therefore gives
\[
 n^2-3\geq n^2-\frac{n^2}{2}=\frac{n^2}{2}>0.
\]
Consequently
\[
 \abs{b_n-1}
 =\frac6{n^2-3}
 \leq\frac6{n^2/2}
 =\frac{12}{n^2}
 \leq\frac{12}{n}.
\]
The last inequality follows from $n^2\geq n>0$ for $n\geq1$.
Fact 1 with $C=12$ and starting index $3$ proves convergence.
The negative first denominator does not enter the tail estimate.

\answer{Part (b) converges to $1$.}

\step{(c) The sequence $c_n=2^{-n}$.}
The expression means the positive reciprocal $1/2^n$.
Fact 2 gives
\[
 \abs{c_n-0}=\frac1{2^n}\leq\frac1{n+1}\leq\frac1n
 \qquad(n\geq1).
\]
The error is thus bounded by $1/n$. Applying Fact 1 proves the
limit directly.

\answer{Part (c) converges to $0$.}

\step{(d) The sequence $t_n=1+2/n$.}
Since $n>0$, every term is defined, and subtracting $1$ gives
\[
 \abs{t_n-1}
 =\abs{1+\frac2n-1}
 =\abs{\frac2n}
 =\frac2n.
\]
For each $\varepsilon>0$, choosing $N>2/\varepsilon$ makes this
error smaller than $\varepsilon$ whenever $n>N$. This is Fact 1
with $C=2$.

\answer{Part (d) converges to $1$.}

\step{(e) The sequence $x_n=73+(-1)^n$.}
If $n$ is even, $(-1)^n=1$ and $x_n=73+1=74$.
If $n$ is odd, $(-1)^n=-1$ and $x_n=73-1=72$.
These values do not merely occur at the start: for any threshold
$N$, choose a positive integer $m>N$. Then $2m>N$ and
$2m+1>N$, and
\[
 x_{2m}=74,\qquad x_{2m+1}=72,\qquad
 \abs{x_{2m}-x_{2m+1}}=\abs{74-72}=2.
\]
Thus Fact 3 applies with $D=2$. No real number can be approached
by every term, because every tail still contains these two values
at distance $2$.

\answer{Part (e) does not converge.}

\step{(f) The sequence $s_n=2^{1/n}$.}
Write $r_n=2^{1/n}$ for the positive real $n$th root of $2$;
thus $r_n>0$ and $r_n^n=2$. We prove the concrete estimate
\[
 0\leq r_n-1\leq\frac1n.
\]
It will imply convergence by Fact 1 without any theorem about
limits of roots or logarithms.

\textbf{First prove the auxiliary Bernoulli inequality.}
For every $x\geq0$ and every positive integer $m$,
\[
 (1+x)^m\geq1+mx.
\]
Fix $x\geq0$. At $m=1$ the two sides are equal. If the inequality
holds at $m$, multiplying by the positive number $1+x$ gives
\begin{align*}
 (1+x)^{m+1}
 &\geq(1+mx)(1+x)\\
 &=1+x+mx+mx^2\\
 &=1+(m+1)x+mx^2\\
 &\geq1+(m+1)x,
\end{align*}
because $mx^2\geq0$. This is the required step at $m+1$,
so induction proves the inequality for the fixed $x$, and hence
for every $x\geq0$.

Apply it with $x=1/n$ and $m=n$. We obtain
\[
 \left(1+\frac1n\right)^n
 \geq1+n\frac1n=2.
\]
For positive real numbers, raising to a positive integer power
preserves strict order. For example, if $u>v>0$, then
$u/v>1$, so the product of $n$ factors $u/v$ is greater than
$1$; multiplication by $v^n>0$ gives $u^n>v^n$.
If $r_n>1+1/n$, this order property would give
\[
 2=r_n^n>\left(1+\frac1n\right)^n\geq2,
\]
a contradiction. Hence $r_n\leq1+1/n$. Also $r_n\geq1$:
if $0<r_n<1$, multiplying $n$ factors each less than $1$
would give $r_n^n<1$, contradicting $r_n^n=2$.

Thus $1\leq r_n\leq1+1/n$, and subtraction of $1$ yields
\[
 \abs{r_n-1}=r_n-1\leq\frac1n.
\]
Fact 1 proves the asserted limit.

\answer{Part (f) converges to $1$.}

\step{(g) The sequence $y_n=n!$.}
For $n\geq1$, the factorial can be written
\[
 n!=n\,(1\cdot2\cdots(n-1)).
\]
At $n=1$, the parenthesized product is empty and equals $1$.
For $n\geq2$, every factor in that product is at least $1$,
so its product is at least $1$. In all cases, therefore,
\[
 n!\geq n.
\]
Now let a real number $B$ and a threshold $N$ be arbitrary.
Choose $n\in\N$ with $n>\max\{N,B,1\}$. Then $n>N$ and
$y_n=n!\geq n>B$. Thus arbitrarily late values exceed every
prescribed real size. Fact 4 rules out convergence to any real
number. This argument establishes the needed nonconvergence
without citing a boundedness theorem.

\answer{Part (g) does not converge to a real limit; its values
grow beyond every real bound.}

\step{(h) The sequence $d_n=(-1)^n n$.}
At an even index $2m$ the value is $2m$, while at an odd index
$2m+1$ the value is $-(2m+1)$. Given any threshold $N$, choose
a positive integer $m>N$. Both indices $2m$ and $2m+1$ are then
greater than $N$, and
\begin{align*}
 \abs{d_{2m}-d_{2m+1}}
 &=\abs{2m-(-(2m+1))}\\
 &=\abs{4m+1}\\
 &=4m+1\geq5>1.
\end{align*}
Here $m\geq1$, so $4m+1$ is positive, justifying the removal
of the absolute value. Every tail has two values separated by
at least $1$. Fact 3 with $D=1$ rules out a real limit.

\answer{Part (h) does not converge. The signs alternate while
the magnitudes increase.}

\step{(i) The sequence $u_n=(-1)^n/n$.}
Although its sign alternates, its distance from $0$ is
\[
 \abs{u_n-0}
 =\frac{\abs{(-1)^n}}{\abs n}
 =\frac1n,
\]
because $\abs{(-1)^n}=1$ and $n>0$. Fact 1 applies with
$C=1$. In contrast with (h), the alternating values now have
small magnitude for every sufficiently large index.

\answer{Part (i) converges to $0$.}

\step{(j) The sequence $u_n=(7n^3+8n)/(2n^3-3)$.}
At $n=1$, the denominator is $2-3=-1$, so it is nonzero.
For $n\geq2$, $n^3\geq8$, giving
$2n^3-3\geq16-3=13>0$. Thus the sequence is defined at every
positive integer.

We test the candidate $7/2$ by exact subtraction:
\begin{align*}
 u_n-\frac72
 &=\frac{2(7n^3+8n)-7(2n^3-3)}{2(2n^3-3)}\\
 &=\frac{14n^3+16n-14n^3+21}{2(2n^3-3)}\\
 &=\frac{16n+21}{2(2n^3-3)}.
\end{align*}
For $n\geq2$, we have $n^3\geq8>3$, so
\[
 2n^3-3\geq n^3>0.
\]
Also $n\geq1$ implies $21\leq21n$ and hence
$16n+21\leq37n$. Using these positive denominator bounds gives
\[
 \abs{u_n-\frac72}
 =\frac{16n+21}{2(2n^3-3)}
 \leq\frac{37n}{2n^3}
 =\frac{37}{2n^2}
 \leq\frac{37}{2n}.
\]
The last step again uses $n^2\geq n>0$. Fact 1, with
$C=37/2$ and starting index $2$, proves convergence.

\answer{Part (j) converges to $7/2$.}

\step{(k) The sequence $u_n=(9n^2-18)/(6n+18)$.}
The denominator $6n+18$ is positive for every $n\geq1$.
We establish a lower estimate that becomes arbitrarily large.
For $n\geq3$, $n^2\geq9$, so $3n^2\geq27>18$ and
\[
 9n^2-18\geq9n^2-3n^2=6n^2>0.
\]
Also $n\geq3$ implies $18\leq6n$, giving
\[
 0<6n+18\leq12n.
\]
First lower the numerator while keeping the positive denominator
fixed, then increase that denominator:
\[
 u_n=\frac{9n^2-18}{6n+18}
 \geq\frac{6n^2}{6n+18}
 \geq\frac{6n^2}{12n}
 =\frac n2.
\]
The second comparison follows from order reversal for positive
reciprocals, followed by multiplication by $6n^2>0$.

Let $B\in\R$ and a threshold $N$ be arbitrary. Choose a positive
integer $n>\max\{N,3,2B\}$. The estimate applies and gives
$u_n\geq n/2>B$, with $n>N$. Fact 4 therefore rules out
convergence to every real candidate.

\answer{Part (k) does not converge to a real limit; its values
grow beyond every real bound.}

\step{(l) The sequence $u_n=\sin(n\pi/2)$.}
Every positive integer has the form $n=4q+r$, with $q\geq0$
an integer and $r\in\{1,2,3,4\}$. Since
$n\pi/2=2q\pi+r\pi/2$, the period $2\pi$ of sine gives
the repeating values
\[
\begin{array}{c|cccc}
 r&1&2&3&4\\ \hline
 \sin(r\pi/2)&1&0&-1&0
\end{array}
\]
The repetition alone suggests nonconvergence; we now check the
arbitrarily late separation required by Fact 3. Given a threshold
$N$, choose $m\in\N$ with $m>N$. Then
$p=4m+1>N$ and $q=4m+3>N$, and periodicity gives
\[
 u_p=\sin(2m\pi+\pi/2)=1,
 \qquad
 u_q=\sin(2m\pi+3\pi/2)=-1.
\]
Their distance is $\abs{1-(-1)}=2$. Fact 3 applies with
$D=2$ and proves nonconvergence.

\answer{Part (l) does not converge.}

\step{(m) The sequence $u_n=\sin(n\pi)$.}
For every integer $n$, the angle $n\pi$ lies at an integer multiple
of a half-turn, so its sine is $0$. Explicitly, if $n=2m$ is
even, then $\sin(n\pi)=\sin(2m\pi)=0$ by periodicity.
If $n=2m+1$ is odd, then
\[
 \sin(n\pi)=\sin(2m\pi+\pi)=\sin\pi=0.
\]
These two cases cover every positive integer. Thus $u_n=0$
for every index, and for every $\varepsilon>0$,
\[
 \abs{u_n-0}=0<\varepsilon
\]
holds at every index. This verifies convergence directly from
the definition.

\answer{Part (m) converges to $0$.}

\step{(n) The sequence $u_n=\sin(2n\pi/3)$.}
Write $n=3q+r$, with $q\geq0$ and $r\in\{1,2,3\}$.
Then $2n\pi/3=2q\pi+2r\pi/3$, giving the period table
\[
\begin{array}{c|ccc}
 r&1&2&3\\ \hline
 \sin(2r\pi/3)&\sqrt3/2&-\sqrt3/2&0
\end{array}
\]
For any threshold $N$, choose a positive integer $m>N$.
The indices $p=3m+1$ and $q=3m+2$ both exceed $N$, and
\[
 u_p=\frac{\sqrt3}{2},\qquad
 u_q=-\frac{\sqrt3}{2}.
\]
Their distance is
\[
 \abs{u_p-u_q}
 =\abs{\frac{\sqrt3}{2}+\frac{\sqrt3}{2}}
 =\sqrt3>0.
\]
Consequently Fact 3 with $D=\sqrt3$ excludes every real limit.

\answer{Part (n) does not converge.}

\step{(o) The sequence $u_n=(\sin n)/n$.}
Here the argument is $n$ radians, with no factor of $\pi$.
We do not need to classify the behavior of $\sin n$ itself.
The elementary trigonometric inequality
$\abs{\sin x}\leq1$ for every real $x$ gives
\[
 \abs{u_n-0}
 =\abs{\frac{\sin n}{n}}
 =\frac{\abs{\sin n}}{n}
 \leq\frac1n,
\]
where division by $n>0$ preserves the inequality.
Fact 1 with $C=1$ proves the limit.

\answer{Part (o) converges to $0$.}

\step{(p) The sequence $u_n=(2^{n+1}+5)/(2^n-7)$.}
The denominator is $2^n-7$. At $n=1$ it
is $2-7=-5$, and at $n=2$ it is $4-7=-3$.
For $n\geq3$, $2^n\geq8$ and therefore $2^n-7\geq1>0$.
These cases show that every term is defined. For example the
first three terms are
\[
 u_1=\frac{4+5}{2-7}=-\frac95,
 \qquad u_2=\frac{8+5}{4-7}=-\frac{13}{3},
 \qquad u_3=\frac{16+5}{8-7}=21.
\]
They do not by themselves determine the limiting behavior.

Subtract the candidate $2$ exactly:
\begin{align*}
 u_n-2
 &=\frac{2^{n+1}+5-2(2^n-7)}{2^n-7}\\
 &=\frac{2\cdot2^n+5-2\cdot2^n+14}{2^n-7}\\
 &=\frac{19}{2^n-7}.
\end{align*}
For $n\geq4$, $2^n\geq16\geq14$, so $7\leq2^n/2$.
It follows that
\[
 2^n-7\geq2^n-\frac{2^n}{2}=\frac{2^n}{2}>0.
\]
Using this bound and Fact 2 yields
\begin{align*}
 \abs{u_n-2}
 &=\frac{19}{2^n-7}\\
 &\leq\frac{19}{2^n/2}
 =\frac{38}{2^n}
 \leq\frac{38}{n+1}
 \leq\frac{38}{n}.
\end{align*}
Fact 1 with $C=38$ and starting index $4$ proves convergence.
It also explains why the negative initial terms and the relatively
large value $u_3=21$ do not change the tail conclusion.

\answer{Part (p) converges to $2$.}

\step{(q) The sequence $a_n=3^n/n!$.}
Both $3^n$ and $n!$ are positive for every $n\in\N$, so
$a_n>0$ and every ratio below is defined. Compare successive
terms using $(n+1)!=(n+1)n!$:
\begin{align*}
 \frac{a_{n+1}}{a_n}
 &=\frac{3^{n+1}}{(n+1)!}\cdot\frac{n!}{3^n}\\
 &=\frac{3\cdot3^n\cdot n!}{(n+1)n!\,3^n}\\
 &=\frac3{n+1}.
\end{align*}
For $n\geq5$, $n+1\geq6$, giving
\[
 0<\frac3{n+1}\leq\frac36=\frac12,
 \qquad
 a_{n+1}\leq\frac12a_n.
\]
We will turn this ratio estimate into an explicit reciprocal bound;
no theorem about ratios of sequences is being invoked.

The value at the starting index is
\[
 C=a_5=\frac{3^5}{5!}
 =\frac{243}{120}=\frac{81}{40}>0.
\]
We claim, for every integer $m\geq0$, that
\[
 a_{5+m}\leq C\,2^{-m}.
\]
At $m=0$, this is equality $a_5=C$. If the estimate holds
at $m$, the ratio bound applies at $n=5+m\geq5$, and gives
\[
 a_{5+(m+1)}=a_{6+m}
 \leq\frac12a_{5+m}
 \leq\frac12C\,2^{-m}
 =C\,2^{-(m+1)}.
\]
This is the next estimate, so induction proves the claim.
Putting $m=n-5$ now yields, for every $n\geq5$,
\[
 0<a_n\leq C\,2^{-(n-5)}.
\]
Fact 2, applied to $m=n-5\geq0$, supplies
\[
 2^{n-5}\geq(n-5)+1=n-4>0.
\]
Taking reciprocals and multiplying by $C>0$ therefore gives
\[
 \abs{a_n-0}=a_n
 \leq C\,2^{-(n-5)}
 \leq\frac{C}{n-4}
 =\frac{81/40}{n-4}.
\]
For any $\varepsilon>0$, choose an integer
$N>\max\{5,4+C/\varepsilon\}$. If $n>N$, then
$n-4>C/\varepsilon$, so the displayed error is less than
$\varepsilon$. This is the shifted estimate of Fact 1, proving
the limit directly.

\answer{Part (q) converges to $0$.}

\step{(r) The sequence $u_n=(1+1/n)^2$.}
Expand by multiplying its two factors, not by assuming a rule
about limits of squares:
\begin{align*}
 u_n
 &=\left(1+\frac1n\right)\left(1+\frac1n\right)\\
 &=1+\frac1n+\frac1n+\frac1{n^2}\\
 &=1+\frac2n+\frac1{n^2}.
\end{align*}
All terms after the leading $1$ are positive. Since
$n^2\geq n>0$, we have $1/n^2\leq1/n$, giving
\[
 \abs{u_n-1}
 =\frac2n+\frac1{n^2}
 \leq\frac2n+\frac1n
 =\frac3n.
\]
Fact 1 with $C=3$ proves convergence.

\answer{Part (r) converges to $1$.}

\step{(s) The sequence $u_n=(4n^2+3)/(3n^2-2)$.}
For all $n\geq1$, $n^2\geq1$, so
$3n^2-2\geq3-2=1>0$. Thus every term is defined and
the denominator has a fixed positive sign.
Subtracting $4/3$ with a common denominator gives
\begin{align*}
 u_n-\frac43
 &=\frac{3(4n^2+3)-4(3n^2-2)}{3(3n^2-2)}\\
 &=\frac{12n^2+9-12n^2+8}{3(3n^2-2)}\\
 &=\frac{17}{3(3n^2-2)}.
\end{align*}
The estimate $n^2\geq1$ also yields
\[
 (3n^2-2)-n^2=2n^2-2=2(n^2-1)\geq0,
\]
so $3n^2-2\geq n^2>0$. Therefore
\[
 \abs{u_n-\frac43}
 =\frac{17}{3(3n^2-2)}
 \leq\frac{17}{3n^2}
 \leq\frac{17}{3n}.
\]
Fact 1 with $C=17/3$ verifies the limit from the definition.

\answer{Part (s) converges to $4/3$.}

\step{(t) The sequence $u_n=(6n+4)/(9n^2+7)$.}
For $n\geq1$, both $6n+4$ and $9n^2+7$ are positive,
so the absolute value of $u_n$ is $u_n$ itself.
We have
\[
 9n^2+7\geq9n^2>0,
 \qquad 6n+4\leq6n+4n=10n,
\]
where $4\leq4n$ uses $n\geq1$. Taking reciprocals in the
denominator comparison and then applying the numerator comparison,
\[
 \abs{u_n-0}
 =\frac{6n+4}{9n^2+7}
 \leq\frac{6n+4}{9n^2}
 \leq\frac{10n}{9n^2}
 =\frac{10}{9n}.
\]
Fact 1 with $C=10/9$ proves convergence to the proposed value $0$.

\answer{Part (t) converges to $0$.}

\subsubsection*{Complete answer list}

In this list, ``does not converge'' means that there is no finite
real limit. Every entry is justified in its corresponding part above.
The growth cases (g) and (k), the persistent oscillations (e), (h),
(l), and (n), and the shrinking oscillations (i) and (o) have been
distinguished by direct estimates and the definition.

\begin{center}
\begin{tabular}{@{}cl@{\qquad}cl@{}}
\toprule
Part&Answer&Part&Answer\\ \midrule
(a)&Converges to $1$&(k)&Does not converge\\
(b)&Converges to $1$&(l)&Does not converge\\
(c)&Converges to $0$&(m)&Converges to $0$\\
(d)&Converges to $1$&(n)&Does not converge\\
(e)&Does not converge&(o)&Converges to $0$\\
(f)&Converges to $1$&(p)&Converges to $2$\\
(g)&Does not converge&(q)&Converges to $0$\\
(h)&Does not converge&(r)&Converges to $1$\\
(i)&Converges to $0$&(s)&Converges to $4/3$\\
(j)&Converges to $7/2$&(t)&Converges to $0$\\
\bottomrule
\end{tabular}
\end{center}
